\section{Dẫn nhập}
\label{sec:introduction}

Dự báo chuỗi thời gian dài hạn (LTSF) từng bị thống trị bởi các kiến trúc kiểu Transformer.
Zeng và cộng sự đã thách thức xu hướng đó bằng \textbf{LTSF-Linear}: họ đề xuất họ baseline
tuyến tính mạnh gồm \textbf{Linear}, \textbf{NLinear} và \textbf{DLinear}, thường ngang
hoặc vượt Transformer trên nhiều benchmark LTSF trong khi vẫn đơn giản và tối ưu chi phí hơn khi huấn luyện
\cite{zeng2023ltsf}.

Ba mô hình dùng cùng kiểu ánh xạ tuyến tính độc lập theo kênh từ độ dài lookback \(L\) sang
horizon \(T\), nhưng khác nhau ở cách xử lý tính không dừng và cấu trúc chuỗi:

\begin{itemize}
      \item \textbf{Linear:} hồi quy trực tiếp từ cửa sổ \(L\) sang \(T\).
      \item \textbf{NLinear (chuẩn hóa):} trừ giá trị lookback cuối, dự báo phần dư, rồi cộng lại:
            đây là cách tái cân bằng (recentering) để chống level shift.
      \item \textbf{DLinear (phân rã):} tách cửa sổ thành thành phần xu hướng (trend) và thành phần mang
            tính mùa vụ (seasonal) bằng nhân trung bình trượt, dự báo mỗi nhánh bằng một đầu Linear riêng rồi
            cộng lại, kế thừa ý tưởng phân rã xu hướng-mùa trong dự báo cổ điển
            \cite{cleveland1990stl,wu2021autoformer,zhou2022fedformer}.
\end{itemize}

Chuẩn hóa và phân rã nhằm các chế độ khác nhau. Chuẩn hóa theo instance (ví dụ RevIN)
hướng tới dịch phân phối giữa train và test \cite{kim2022revin}. Phân rã nhắm cấu trúc xu hướng
và mùa vụ rõ ràng. Trên benchmark "sạch", Toner và Darlow cho thấy DLinear thường gần như
cùng lớp mô hình với Linear thuần \cite{toner2024linear}, nhưng điều đó vẫn \emph{không} trả lời
một người thực hành: \emph{khi nào} nên chọn DLinear thay vì NLinear.

Do đó, dự án tập trung vào một câu hỏi so sánh duy nhất:
\textbf{Khi nào Decomposition vượt qua Normalization?}
Chúng ta xem DLinear và NLinear như hai inductive bias, và hỏi đặc trưng dữ liệu nào dự đoán
mô hình nào thắng, thay vì chỉ báo mô hình nào có MSE thấp hơn trên từng ô bảng xếp hạng.
